\section*{Sets and Indices}

There are no nontrivial sets or indices in this model. The formulation has two continuous production decision variables, one for each product:
\begin{itemize}
    \item Product A,
    \item Product B.
\end{itemize}

\section*{Parameters}

\begin{itemize}
    \item $p_A$ (code: \texttt{profit\_A}): profit per kilogram of product A.
    \item $p_B$ (code: \texttt{profit\_B}): profit per kilogram of product B.
    \item $H$ (code: \texttt{max\_hours}): maximum production hours available per week.
    \item $h_A$ (code: \texttt{hours\_A}): production hours required per kilogram of product A.
    \item $h_B$ (code: \texttt{hours\_B}): production hours required per kilogram of product B.
    \item $r^{BA}$ (code: \texttt{min\_ratio\_B\_to\_A}): minimum output ratio requiring product B production to be at least $r^{BA}$ times product A production.
    \item $s^{AB}$ (code: \texttt{storage\_ratio\_A\_to\_B}): storage-space ratio, where one kilogram of product A uses $s^{AB}$ times the storage space of one kilogram of product B.
    \item $\bar{x}_A^{\text{stor}}$ (code: \texttt{max\_storage\_A}): storage capacity expressed as an equivalent upper bound of product A alone.
    \item $\bar{S}$ (code: \texttt{max\_storage\_units}): total storage capacity in storage-space units, computed in the code as
    \[
    \bar{S} = \bar{x}_A^{\text{stor}} \, s^{AB}.
    \]
\end{itemize}

For this instance, the data imply
\[
p_A=30,\quad p_B=10,\quad H=40,\quad h_A=6,\quad h_B=3,\quad r^{BA}=3,\quad s^{AB}=4,\quad \bar{x}_A^{\text{stor}}=4,
\]
and therefore
\[
\bar{S}=4\cdot 4=16.
\]

\section*{Decision Variables}

\begin{itemize}
    \item $x_A$ (code: \texttt{x\_A}): kilograms of product A to produce.
    \item $x_B$ (code: \texttt{x\_B}): kilograms of product B to produce.
\end{itemize}

Both variables are continuous and nonnegative:
\[
x_A \ge 0,\qquad x_B \ge 0.
\]

\section*{Objective}

Maximize total weekly profit:
\[
\max \; p_A x_A + p_B x_B.
\]

For this instance:
\[
\max \; 30x_A + 10x_B.
\]

\section*{Constraints}

\begin{align}
h_A x_A + h_B x_B &\le H
&& \text{(production hours)} \\
x_B &\ge r^{BA} x_A
&& \text{(market demand ratio)} \\
s^{AB} x_A + x_B &\le \bar{S}
&& \text{(storage capacity)} \\
x_A,\; x_B &\ge 0.
\end{align}

Substituting the instance values gives:
\begin{align}
6x_A + 3x_B &\le 40, \\
x_B &\ge 3x_A, \\
4x_A + x_B &\le 16, \\
x_A,\; x_B &\ge 0.
\end{align}

\section*{Notes on Implementation}

The source code builds this model as a linear program with two continuous variables using PuLP-compatible syntax and solves it with \texttt{GUROBI\_CMD}.

The storage constraint is implemented exactly as
\[
\texttt{storage\_ratio\_A\_to\_B} \cdot x_A + x_B \le \texttt{max\_storage\_A} \cdot \texttt{storage\_ratio\_A\_to\_B},
\]
i.e.,
\[
s^{AB} x_A + x_B \le \bar{x}_A^{\text{stor}} s^{AB}.
\]
Thus, \texttt{max\_storage\_A} is not used as a direct bound on $x_A$; instead, it is converted into total storage-space units before defining the single storage-capacity inequality.

There are no integer restrictions, no additional bounds, and no hidden reformulations beyond the linear constraints shown above.
